\documentclass[11pt]{article}

\usepackage[a4paper,margin=3cm]{geometry}
\usepackage{microtype}
\usepackage{amsmath,amssymb,amsthm}
\usepackage{xcolor}
\usepackage{url}

% Every count below is generated from the run stores by tools/note_tables.py and
% tools/density.py. Prose that spells a number out goes stale the moment the next
% area lands, and goes stale quietly, because nothing rereads a paragraph.
\newcommand{\mathlibTheorems}{186{,}856}
\newcommand{\mathlibExaminedPercent}{15.7}
\newcommand{\mathlibWeakenablePercent}{0.23}
\newcommand{\reweightedRate}{1.71}
\newcommand{\reweightAreas}{15}
\newcommand{\reachablePercent}{67}
\newcommand{\noTableClassPercent}{17}
\newcommand{\noBinderPercent}{16}
\newcommand{\tauTheorems}{29{,}589}
\newcommand{\tauExamined}{5{,}526}
\newcommand{\tauWeakenable}{62}
\newcommand{\tauExaminedPercent}{18.7}
\newcommand{\tauOfExaminedPercent}{1.12}
\newcommand{\tauWeakenablePercent}{0.21}
\newcommand{\citedN}{571}
\newcommand{\citedNever}{174}
\newcommand{\citedNeverPercent}{30.5}
\newcommand{\citedNeverCI}{[26.3, 35.3]}
\newcommand{\citedDeff}{1.42}
\newcommand{\citedControlPercent}{37.7}
\newcommand{\citedControlN}{20{,}125}
\newcommand{\citedGap}{7.2}
\newcommand{\citedP}{0.0018}
\newcommand{\citedMean}{2.34}
\newcommand{\citedMax}{193}
\newcommand{\citedControlMean}{4.28}
\newcommand{\citedControlMax}{4{,}249}
\newcommand{\citedGapClean}{12.4}
\newcommand{\citedGapAll}{15.7}
\newcommand{\citedGenerated}{22{,}854}
\newcommand{\patchGroups}{102}
\newcommand{\patchTested}{68}
\newcommand{\patchKept}{36}
\newcommand{\patchRejected}{32}
\newcommand{\namedLocated}{293}
\newcommand{\namedDocumented}{47}
\newcommand{\namedClassical}{1}
\newcommand{\namedDocumentedPercent}{16}
\newcommand{\unusedCount}{541}
\newcommand{\replaceableCount}{77}
\newcommand{\unclassifiedCount}{27}
\newcommand{\classifiedCount}{618}
\newcommand{\unusedPercent}{87.5}
\newcommand{\replaceablePercent}{12.5}
\newcommand{\areasSwept}{4}
\newcommand{\areasSweptWord}{four}
\newcommand{\areaNames}{Algebra, Analysis, NumberTheory, RingTheory}
\newcommand{\attempts}{73{,}620}
\newcommand{\survivors}{322}
\newcommand{\lostInvertibility}{44}
\newcommand{\deepest}{0}
\newcommand{\deepestLost}{0}
\newcommand{\libraryFiles}{8{,}311}
\newcommand{\sweptFiles}{3{,}105}
\newcommand{\sweptPercent}{37}
\newcommand{\filterRejectPercent}{52}
\newcommand{\ourYield}{0.44}
\newcommand{\ourYieldCI}{[0.39, 0.49]}
\newcommand{\incoherentPercent}{75}
\newcommand{\tableEdges}{48}
\newcommand{\attemptsPerFile}{24}
\newcommand{\decisionsPerMinute}{26}
\newcommand{\freeVerdictPercent}{58}
\newcommand{\bestTotal}{2{,}877}
\newcommand{\bestTop}{25}
\newcommand{\bestMissing}{13}
\newcommand{\bestHeadCount}{96}
\newcommand{\bestHeadAttempts}{2{,}922}
\newcommand{\bestHeadFound}{3}
\newcommand{\bestHeadNotABinder}{81}
\newcommand{\bestHeadIncoherent}{73}
\newcommand{\bestFieldDR}{23}
\newcommand{\ourFieldDR}{25}
\newcommand{\predictedAreas}{MeasureTheory, Order, Topology}
\newcommand{\predictedAttempts}{13{,}525}
\newcommand{\predictedSurvivors}{67}
\newcommand{\predictedDeep}{0}
\newcommand{\predictedDeepExpected}{0.0}
\newcommand{\predictedDeepP}{1.000}
\newcommand{\baseLostPercent}{13.7}
\newcommand{\predictedLost}{3}
\newcommand{\predictedLostPercent}{4.5}
\newcommand{\predictedLostExpected}{9.2}
\newcommand{\predictedLostP}{0.0383}
\newcommand{\allAreas}{18}
\newcommand{\allAreasWord}{18}
\newcommand{\allAreasWordCap}{18}
\newcommand{\allAreaNames}{Algebra, AlgebraicGeometry, Analysis, Combinatorics, Data, Dynamics, FieldTheory, Geometry, GroupTheory, LinearAlgebra, MeasureTheory, ModelTheory, NumberTheory, Order, Probability, RepresentationTheory, RingTheory, Topology}
\newcommand{\allSurvivors}{462}
\newcommand{\allAttempts}{116{,}265}
\newcommand{\allSweptFiles}{6{,}413}
\newcommand{\allSweptPercent}{77}
\newcommand{\allYield}{0.40}
\newcommand{\allYieldCI}{[0.36, 0.44]}
\newcommand{\theoremsExamined}{29{,}261}
\newcommand{\theoremsWeakened}{431}
\newcommand{\theoremRate}{1.47}
\newcommand{\theoremOneIn}{67}
\newcommand{\theoremFiles}{240}
\newcommand{\allAttemptsPerFile}{18}
\newcommand{\derivedAreas}{25}
\newcommand{\derivedSweptFiles}{7{,}869}
\newcommand{\derivedSweptPercent}{95}
\newcommand{\derivedAttempts}{148{,}634}
\newcommand{\gradientRows}{LinearAlgebra & 10/22 & 45\% \\ Combinatorics & 2/8 & 25\% \\ GroupTheory & 8/9 & 89\%}
\newcommand{\gradientMax}{89}
\newcommand{\gradientLost}{20}
\newcommand{\gradientTotal}{39}
\newcommand{\gradientPercent}{51}
\newcommand{\gradientP}{3.1\times 10^{-8}}
\newcommand{\rerunRows}{Algebra & 151 & 139 & 89 & 201 \\ Analysis & 107 & 53 & 15 & 145 \\ RingTheory & 53 & 20 & 7 & 66 \\ Order & 30 & 44 & 29 & 45 \\ Topology & 25 & 29 & 12 & 42 \\ LinearAlgebra & 22 & 25 & 9 & 38 \\ MeasureTheory & 12 & 10 & 1 & 21 \\ Data & 13 & 10 & 6 & 17 \\ NumberTheory & 11 & 8 & 5 & 14 \\ GroupTheory & 9 & 14 & 9 & 14 \\ Combinatorics & 8 & 11 & 7 & 12 \\ Geometry & 9 & 7 & 5 & 11 \\ Probability & 5 & 4 & 2 & 7 \\ FieldTheory & 3 & 4 & 2 & 5 \\ Dynamics & 2 & 1 & 1 & 2}
\newcommand{\rerunHand}{460}
\newcommand{\rerunUnion}{640}
\newcommand{\rerunGainPercent}{39}
\newcommand{\catArea}{CategoryTheory}
\newcommand{\catFiles}{1{,}089}
\newcommand{\catAttempts}{703}
\newcommand{\catSurvivors}{1}
\newcommand{\catContextPercent}{28}
\newcommand{\derivedEdges}{374}
\newcommand{\derivedClasses}{224}
\newcommand{\denseArea}{LinearAlgebra}
\newcommand{\denseSurvivors}{22}
\newcommand{\denseLost}{10}
\newcommand{\denseLostPercent}{45}
\newcommand{\denseVsPredictedP}{2.3\times 10^{-5}}
\newcommand{\denseVsBaseP}{0.0006}
\newcommand{\tauVsMathlibP}{0.35}
\newcommand{\tauRatio}{1.13}
\newcommand{\tauAttempts}{15{,}602}
\newcommand{\tauSurvivors}{70}
\newcommand{\tauYield}{0.45}
\newcommand{\tauYieldCI}{[0.36, 0.57]}

\newcommand{\densityAreas}{21}
\newcommand{\densityPooled}{9.8}
\newcommand{\densityG}{126}
\newcommand{\densityDf}{20}
\newcommand{\densityP}{2.7e-17}
\newcommand{\densityDeff}{2.6}
\newcommand{\densityGClustered}{49}
\newcommand{\densityPClustered}{0.0003}
\newcommand{\sweptAreas}{25}
\newcommand{\sweptFilesAll}{7{,}869}
\newcommand{\sweptPercentAll}{95}
\newcommand{\unsweptFiles}{442}
\newcommand{\unsweptAreas}{Tactic, Lean, Util, InformationTheory, Testing, Deprecated}
\newcommand{\densityDensest}{Combinatorics}
\newcommand{\densityDensestRate}{29.9}
\newcommand{\densityDensestDecls}{368}
\newcommand{\densitySparsest}{RingTheory}
\newcommand{\densitySparsestRate}{3.9}
\newcommand{\densitySparsestDecls}{5151}
\newcommand{\densitySpread}{8}
\newcommand{\densityCrossRho}{+0.80}
\newcommand{\densityCrossN}{20}
\newcommand{\densityCrossP}{2e-05}


\newtheorem{observation}{Observation}

% No subtitle. It restated the abstract's opening sentence, and it named Mathlib
% where that sentence says "a large formal library" -- which is the phrasing that
% covers the machine-written corpus in Section 5 as well.
\title{Assumptions that were never used\thanks{Produced with AI assistance throughout --- tooling, sweep, statistics and drafting --- for which the author is responsible. No claim here rests on a model's judgement: every row is a Lean compile with its axioms checked and every figure is generated from the published data.}}
\author{Carlo Perassi}
\date{September 2026}

\begin{document}
\maketitle



\begin{abstract}
We describe a mechanical search for theorems in a large formal library whose
stated algebraic setting is stronger than the setting their own proof requires.
The question is not new: it was answered analytically for this library in 2021,
and the corresponding brute-force method --- weaken the hypothesis, re-run the
proof --- was carried out on a different library in 2013 and judged, in the same
sentence that cites it, too slow for this one. We report what it costs, and we
run it over two libraries rather than one.

We examined a \emph{sample} of \theoremsExamined{} theorems of Mathlib, across
\allAreasWord{} subject areas, and found \theoremsWeakened{} that hold with a
weaker hypothesis than the one they are stated with: \theoremRate{}\%, about one
in \theoremOneIn{}, spread over \theoremFiles{} files. The sample is not random --- it is what the search could reach --- and
re-weighting it to the library's own shape gives \reweightedRate{}\%, so the
composition costs little and costs it in our favour. The same instrument on Tau Ceti
\cite{tauceti} --- a corpus of formal mathematics written by AI contributors
rather than by people --- finds them at \tauYield{}\% of attempts against
\allYield{}\% here, intervals \tauYieldCI{} and \allYieldCI{}, $p =
\tauVsMathlibP{}$. We can measure no difference at a resolution that would have
shown one of half again as much. Whatever this is, it is not a habit of one
library's authors.

Per declaration the method remains some fifty times more expensive than reading
the proof term, so it is not the interactive tool that judgement was about; but
a library fits in a few days of one machine, which is a bound nobody had put on
it. None of the survivors is new mathematics, and in \unusedPercent{}\% of the
rows we could check the weakened proof is the same proof term, so the assumption
was genuinely unused rather than merely replaceable by a search. What they give
up is invertibility --- subtraction or division --- and in Mathlib they do so
far more often in
areas built on invertible structure than in areas without it. The phenomenon is
not spread evenly across subjects either: one rate does not explain the areas we
swept, the extremes differ by a factor of \densitySpread{}, and we do not know
why.
\end{abstract}

\section{What we did and what came out}

Take a theorem from Mathlib. Replace one instance binder with a weaker class.
Keep the proof byte for byte. Compile the result alone, and keep it only if Lean
reports no error and \texttt{\#print axioms} shows nothing beyond
\texttt{propext}, \texttt{Classical.choice} and \texttt{Quot.sound}. No
judgement enters, and nothing survives that the compiler has not re-checked.

Three things came out, and they are the whole of what we claim.

\paragraph{One. It is affordable.} A pass over the library is about four days on
one laptop, covering \derivedSweptFiles{} of \libraryFiles{} source files,
\derivedSweptPercent{}\%. That contradicts nothing
anyone measured, and settles something nobody had: the method is a batch
instrument, not an editor-time one, and its batch cost is bounded.

\paragraph{Two. What gets given up is invertibility, and where depends on the
area.} Of the theorems we can weaken in the \areasSweptWord{} areas we swept
first,
\lostInvertibility{} --- \baseLostPercent{}\% --- remove the ability to subtract
or divide. In areas built on invertible structure the rate is
\gradientPercent{}\% ($p = \gradientP{}$); in areas with nothing to invert it is
\predictedLostPercent{}\%. We pre-registered that contrast and tested it on seven
areas that played no part in forming it: $25\%$ against $0\%$, $p = 0.016$.

\paragraph{Three. The phenomenon is not spread evenly.} Across \densityAreas{}
completed areas the rate runs from nothing to \densityDensestRate{} per thousand
declarations. One rate does not explain them: $G = \densityGClustered{}$ on
\densityDf{} degrees of freedom, $p = \densityPClustered{}$, clustering by file.

One theorem can be weakened in more than one way, so \theoremsWeakened{}
theorems give \allSurvivors{} distinct weakenings; the method counts attempts
rather than theorems, and Section~\ref{sec:method} says how the two relate.
A fourth, which we set out to test rather than to find, is the
\unusedPercent{}\% in the abstract: the weakened proof is usually the same proof
term. Sections~\ref{sec:cost} to \ref{sec:density} take these in turn.
Section~\ref{sec:limits} says what none of it shows, which is more than usual.
Four worked examples are in Appendix~\ref{app:obs}; all \allSurvivors{} are in
the companion catalogue, unfiltered, because no filter we have can judge them.

\section{Method}\label{sec:method}

\paragraph{The substitution.} A weakening table maps each class to classes
weaker than it. We use two, and they are not nested in either direction. One is
hand-written, \tableEdges{} edges over the hierarchy as we read it. The other is
derived from the library's own \texttt{class \dots\ extends} declarations,
\derivedEdges{} edges over \derivedClasses{} classes. Over the fifteen areas
swept both ways the hand table finds \rerunHand{}, the derived \(379\), with
\(199\) in common and \rerunUnion{} in the union --- \rerunGainPercent{}\% more
than either alone. A derived table cannot propose a \emph{sibling} edge, and
some of the most productive edges are between siblings; a hand table cannot name
classes nobody thought to list. Every count here is therefore a lower bound, and
the factor is measured rather than guessed.

\paragraph{Three filters before any compile.} Most proposed weakenings are dead
on arrival, and rejecting them cheaply is what makes the pass affordable.
\freeVerdictPercent{}\% of attempts are settled without invoking the compiler:
the binder named is not one the theorem carries, or the assumption never reaches
the conclusion, or the substitution is not a weakening. One \texttt{\#check} per
declaration decides these.

\paragraph{A prior-art gate.} Every nomination is put to \texttt{exact?} on the
weakened statement, and dropped unless nothing usable comes back. A row where
the library's own search closes the goal is not our finding. A row whose
citation is the theorem's own name means the binder was never a hypothesis.

\paragraph{Theorems, weakenings, attempts.} Three units. The paper uses
whichever one the sentence is about, and they differ by more than rounding, so
the table says which is which.

\begin{center}
\begin{tabular}{llr}
\textbf{unit} & \textbf{what it counts} & \textbf{Mathlib}\\[2pt]
theorem examined & a declaration for which at least one weakening & \\
                 & was proposed and decided & \theoremsExamined{}\\[2pt]
theorem weakenable & of those, one that holds with a weaker & \\
                   & hypothesis than it is stated with & \theoremsWeakened{}\\[2pt]
weakening & one binder of one theorem holding over one & \\
          & weaker class; a theorem may have several & \allSurvivors{}\\[2pt]
attempt & one candidate class put to the compiler; a & \\
        & theorem generates many & \allAttempts{}\\
\end{tabular}
\end{center}

The two corpora, on both denominators:

\begin{center}
\begin{tabular}{lrr}
 & \textbf{Mathlib} & \textbf{Tau Ceti}\\[2pt]
theorems in the corpus      & \mathlibTheorems{} & \tauTheorems{}\\
examined                    & \theoremsExamined{} & \tauExamined{}\\
\quad as \% of the corpus   & \mathlibExaminedPercent{}\% & \tauExaminedPercent{}\%\\[2pt]
weakenable                  & \theoremsWeakened{} & \tauWeakenable{}\\
\quad as \% of examined     & \theoremRate{}\% & \tauOfExaminedPercent{}\%\\
\quad as \% of the corpus   & \mathlibWeakenablePercent{}\% & \tauWeakenablePercent{}\%\\
\end{tabular}
\end{center}

About one theorem in six is examined at all; the rest carry no instance binder
either table can weaken. So \theoremRate{}\% is the rate among theorems the
method can ask about and \mathlibWeakenablePercent{}\% is the rate in the
library; Section~\ref{sec:limits} says why neither alone is the honest number.
Both corpora agree on both denominators, which is a stronger statement than the
yield comparison of Section~\ref{sec:tauceti} alone: the rates there are over
attempts, and these are over theorems. Attempts are reported for cost, not for
yield.

Every survivor is a \emph{weakening}: a class replaced by a weaker one. No
binder is ever deleted, so nothing here counts an assumption removed outright.

\paragraph{Which set each number counts.} Three populations appear below and are
not interchangeable, so each is named where it is used. The \emph{discovery set}
is \areaNames{}: \survivors{} survivors from \attempts{} attempts, where the
pattern was found, and every statistic about the pattern is computed over it.
The \emph{full hand-table sweep} is \allAreasWord{} areas, \allSurvivors{}
survivors from \allAttempts{} attempts, and contains the discovery set. The
\emph{derived-table sweep} is a second pass with the other table, reported
separately throughout; Section~\ref{sec:density} is computed over it and says
so.

\section{The cost}\label{sec:cost}

Best \cite{best2021} inspects a theorem's elaborated proof term and reports
which typeclass fields the proof never touches. In the same paper he sets our
variant aside, and we quote him entire because the elision matters:

\begin{quote}
a brute force strategy of weakening typeclasses in theorem statements and
re-running the same proof script seems far too slow to be used on the scale of a
large library in a prover such as Lean, however similar techniques have been
used in the Mizar system [12]. Nevertheless, a tool to automate this process
could still be useful for small new developments.
\end{quote}

His [12] is Alama \cite{alama2013}, who did exactly this to the Mizar
Mathematical Library a decade before. On the reading Best meant --- a tool
usable while formalising --- he was right, and our own numbers confirm it: term
inspection handles a $475$-declaration file in one to two minutes, and
substituting and recompiling the same file costs us something over an hour.
Fifty times more expensive per declaration, and no hardware since closes that.

What we can add is the bound he did not state. At
\allAttemptsPerFile{} attempts per file and \decisionsPerMinute{} decisions a
minute, a complete pass is a few days of one machine. Not a cluster, not a
grant. The derived table has now been over \derivedAreas{} areas,
\derivedSweptFiles{} files, \derivedSweptPercent{}\% of the library; the
\unsweptFiles{} files left are \unsweptAreas{}, which hold no theorems of the
kind this looks for. The extrapolation is no longer an extrapolation.

\section{Invertibility}\label{sec:gradient}

Of the \survivors{} survivors in the \areasSweptWord{} discovery areas
(\areaNames{}), \lostInvertibility{} --- \baseLostPercent{}\% --- give up
subtraction or division. That is the baseline. The \allSurvivors{} of the
abstract is the full hand-table sweep, \allAreasWord{} areas; every statistic
about the pattern is computed over the \survivors{}, because that is where the
pattern was found.

The claim is about where. An area whose objects carry inverses --- groups,
rings, fields, modules over them --- should yield survivors that surrender them;
an area with no inverses to surrender cannot. We named three areas and predicted
each before sweeping it:

\begin{center}
\begin{tabular}{lrr}
\gradientRows{}
\end{tabular}
\end{center}

Against \baseLostPercent{}\% at baseline and \predictedLostPercent{}\% in areas
where these weak structures are already the subject, the three together give
\gradientLost{} of \gradientTotal{} --- \gradientPercent{}\%,
$p = \gradientP{}$.

\emph{One of the three predictions was wrong.} We expected Combinatorics low,
reasoning that counting arguments run over the naturals where subtraction is
unavailable. It sits high, and the explanation --- that this part of the library
is largely additive combinatorics, conducted in groups --- is one we can give
only after seeing the number. We report it because a scale reported only where
it worked is selection rather than evidence.

\paragraph{The statistic was chosen after seeing the data, so we re-ran it as a
pre-registration.} The hypothesis, the classification of each area, the test and
the significance level were written down and committed before the measurement.
Seven areas that had played no part in forming the claim: those built on
invertible structure gave $4$ of $16$, those without gave $0$ of $26$. Fisher
exact, two-sided, $p = 0.016$. The gradient is much shallower than the post-hoc
figure --- $25\%$ against $\gradientPercent{}\%$ --- and what carries it is the
low end, $0$ of $26$ in areas with nothing to invert. Read the
\gradientPercent{}\% as what post-hoc selection does to a real but modest effect.

\section{Never used, or merely not needed?}\label{sec:unused}

A row says the theorem compiles with a weaker hypothesis. That covers two facts.
Either the elaborated proof term is the same one and the assumption was dead
weight, or the tactic, re-run in the weaker setting, searched again and found
another route --- in which case the assumption was not unused but
\emph{replaceable}. Best predicted the second and did not measure it, expecting
it ``especially likely with small proofs that make heavy use of automation''.

We measure it by elaborating each weakened declaration alongside the original,
reading both proof terms out of the environment, and comparing the constants
each cites. Typeclass plumbing is excluded --- substituting a class necessarily
changes which classes, parent projections and instances a term mentions, so
counting those would call every row replaceable. What remains is the lemmas the
proof appeals to.

Of \classifiedCount{} rows we could classify, \unusedCount{} ---
\unusedPercent{}\% --- cite exactly the same lemmas: the same proof, and the
assumption was never used. \replaceableCount{} (\replaceablePercent{}\%) cite
different ones. \unclassifiedCount{} could not be classified at all, because the
original declaration is \texttt{private} and cannot be named from outside its
file; those rows verify like the rest and are reported separately rather than
folded in.


A companion project \cite{companion} measured the same question at a different
strength, reading only the \emph{root} of each proof term --- the lemma the proof ends on --- and
found it unchanged in $99.15\%$ of comparable rows against our
\unusedPercent{}\% over the whole cited set. The two are nested, and joining
them row by row leaves the impossible cell empty: no row has a moved root and an
unchanged citation set. That empty cell is the only evidence either project has
that the two measurements measure what they claim; a single row in it would have
convicted one pipeline without saying which. The gap between the two is the
quantity neither can state alone. Roughly one weakening in eight keeps its
destination and changes its route.

Two things this did not show, both predicted. The gradient of
Section~\ref{sec:gradient} is unchanged when restricted to the genuinely unused
rows --- $17.4\%$ of them give up invertibility against $17.4\%$ of all rows,
$p = 0.63$ --- so nothing in that argument rests on the distinction. And
replaceability is not concentrated in proofs that visibly search: $14\%$ of
rows using \texttt{simp} or a similar tactic are replaceable against $12\%$ of
those using none, $p = 0.36$. Whatever makes a proof replaceable, it is not
simply that a tactic went looking.

\section{The same rate in a library people did not write}\label{sec:tauceti}

If the pattern were a habit of Mathlib's authors, a corpus whose proofs people
did not write should not share it. Tau Ceti \cite{tauceti} is the closest thing
available: formal mathematics implemented by AI contributors from human-written
roadmaps, under adversarial review, incubated by the Lean FRO and the Mathlib
Initiative. People set its targets and judge its output but do not write its
proofs, and it is proofs that carry the assumptions we count. Swept with the same hand-written table --- 42 of its 48 edges appear in that
store and none outside it --- it yields \tauSurvivors{} survivors from
\tauAttempts{} attempts, \tauYield{}\%, against \allYield{}\% over the full
hand-table sweep here. Intervals \tauYieldCI{} and \allYieldCI{}; Fisher exact,
two-sided, $p = \tauVsMathlibP{}$; ratio \tauRatio{}. They overlap.

Say what that can and cannot exclude. At \tauAttempts{} attempts against
\allAttempts{} the comparison has power $0.96$ against a corpus twice as
over-constrained and $0.78$ against one half again as much. A difference of
$50\%$ would probably have shown. This is a null result with a stated
resolution, not a demonstration that the rates are equal --- and the resolution
is the point: an earlier pass on a $4{,}000$-attempt sample of the same corpus
could only exclude a doubling.

The comparison is of \emph{yield}. Whether the invertibility gradient of
Section~\ref{sec:gradient} also survives a change of author is a separate
question and we cannot answer it here: Tau Ceti's \(25\) gated survivors sit
almost entirely in areas built on invertible structure, so there is no low group
to compare against. Its survivors give up invertibility at \(6\) of \(25\)
against \(\lostInvertibility{}\) of \(\survivors{}\) here, $p = 0.23$, which is
consistent and settles nothing.

No prior work makes this comparison. Best ran a single pass over one library and
says so; the question of whether unused assumptions are an artefact of human
authorship needs a corpus no human wrote, and one now exists.

Attempts is the wrong denominator for a comparison between libraries and we use
it because it is the only one available: a table with more edges per class
produces more attempts per declaration and so a lower rate at identical yield.
The defensible sentence is that we can measure no difference, not that the rates
are equal.

\section{Distribution}\label{sec:density}

Counting survivors rewards an area for being large. The question underneath is
whether a theorem in one area is \emph{more likely} to carry an unused
assumption, which needs a denominator: declarations for which some weakening was
proposed. The numerator is theorems with at least one survivor, not survivors,
since one theorem can yield three and those are not three independent trials.

This section alone is computed over the derived-table sweep, which reached the
end of the library. Across \densityAreas{} completed areas the rate runs from
nothing to \densityDensestRate{} per thousand declarations, pooled at
\densityPooled{}. Taking theorems as independent trials gives
$G = \densityG{}$, $p = \densityP{}$ --- and they are not independent trials.
Mathlib heads a file with \texttt{variable [Field K]} and every theorem below
inherits it, so the assumption being weakened belongs to the file at least as
much as to the theorem. At the observed design effect of \densityDeff{} the
statistic falls to $G = \densityGClustered{}$,
$p = \densityPClustered{}$, which still rejects one rate. Under the hand table
the same correction leaves no significant result, $p = 0.13$, and we report that
rather than only the instrument that agrees with us.

Two objections, both answerable from the same data. A dense area might simply
have had more attempts per declaration, but the rank correlation between rate
and attempts per declaration is $-0.23$ ($p = 0.31$). And the ranking might be a
picture of our weakening table rather than of the library; re-measured under the
other table the levels move a great deal and the order holds at
$\rho = \densityCrossRho{}$ over \densityCrossN{} areas
($p = \densityCrossP{}$). The level is the instrument's. The order largely is
not.

Two things this is not. The two tables are not independent, both being orderings
of the same hierarchy, so that agreement is agreement rather than replication.
And the denominator counts only declarations some table could propose a
weakening for, so this is density among the theorems the method can ask about,
which is not density among theorems. We do not know what the second would be.

\section{Who uses them}\label{sec:used}

A weakening that holds is not thereby wanted. Naming conventions speak to a
library's habits; citations speak to its use.

A companion project \cite{companion} counted, over the whole library at our
revision, how many declarations cite each survivor --- from proof terms, not from names. Of the
\citedN{} survivors it could join, \citedNever{} are never cited at all:
\citedNeverPercent{}\%, interval \citedNeverCI{} after correcting for clustering
by module at a design effect of \citedDeff{}. Against non-survivors \emph{in the
survivors' own modules} --- \citedControlN{} theorems, same files, same authors,
same density of machinery --- \citedControlPercent{}\% are never cited. The gap
is \citedGap{} points, survivors lower, $p = \citedP{}$.

Wider comparisons give a wider gap, and the difference between them says why:
\citedGapAll{} points against every other theorem in the corpus, and
\citedGapClean{} after dropping the \citedGenerated{} declarations Mathlib
generates rather than writes (\texttt{.injEq}, \texttt{.noConfusion}, match
auxiliaries), which are largely never cited by name. About half the raw
difference is machinery nobody cites and nobody wrote. The same-module control
is the honest comparison and is the one we quote.

So a redundant hypothesis is not a marker of dead code, and not a marker of
infrastructure either. It sits in theorems used a handful of times each and then
not examined again: mean \citedMean{} against \citedControlMean{} for their
neighbours, and the most-cited survivor is cited \citedMax{} times against
\citedControlMax{} for the most-cited neighbour. A lemma a thousand proofs
depend on has been read by many people; a lemma nothing depends on was perhaps
never finished.

Mathlib's own conventions say the same thing more weakly. Of the
\namedLocated{} survivor declarations our locator could find in the source,
\namedDocumented{} carry a docstring at all (\namedDocumentedPercent{}\%), and
\namedClassical{} has a docstring naming a classical result in bold, which is
how Mathlib marks one: the \emph{Shrinking lemma}, weakened from a metric space
to a pseudometric one. A handful of others are recognisable ---
\texttt{riesz\_lemma}, the Hadamard three-lines bounds --- but by surname in the
declaration name only 22 of \allSurvivors{} match, and most of those are
namespace coincidences rather than named theorems. Those tests are about naming
rather than use, and our locator failed on a quarter of the rows. The citation
counts are the measurement and they agree.

\section{What can actually be applied}\label{sec:patches}

A row is not a patch. \(89\%\) of the assumptions we weaken are written on a
\texttt{variable} line, so they belong to a section rather than to a theorem:
weakening one there changes every declaration below it, and most of those still
need the stronger class. A patch set derived row by row would be wrong more
often than right.

There is a subset where it is safe, and we let the compiler decide which.
\patchGroups{} groups have the property that every theorem carrying the binder
also survived --- the property one would call safe on the evidence. Of the
\patchTested{} where the binder is actually on a \texttt{variable} line, we
rewrote it, rebuilt the whole file against the pinned revision, and kept what
compiled: \textbf{\patchKept{}}. The compiler rejected \patchRejected{}, very nearly
half.

That rejection rate is the useful number. A declaration the search never tried
can still need the stronger class, and no amount of reasoning over our own
verdicts finds it --- only building the file does.

The \patchKept{} that survive are one-line changes against a named revision,
shipped in \texttt{patches/} with every rejection and its first compiler error
beside them. They are the smallest actionable thing this work produces, and the
only part of it a maintainer can evaluate without reading the rest.

A patch against a pinned revision decays, so we measured the decay. At
\texttt{\patchStillHoldRev{}}, $804$ commits and a toolchain version later,
\patchStillHold{} of the \patchKept{} still rebuild their whole file and none
applied and then failed. Of the three that did not carry over,
\patchSuperseded{} were superseded upstream in the same direction and one file
was rewritten. Published as \texttt{patches/REPORT-master.json}; no figure
elsewhere depends on it.

\section{Relation to prior work}\label{sec:prior}

The question has been asked of this library and answered analytically. Best
\cite{best2021} produced $2{,}877$ suggested generalisations in a single pass by
reading elaborated proof terms. The substitute-and-recompile method is older and
belongs to Alama \cite{alama2013} on the Mizar Mathematical Library; Pons
\cite{pons2002} is the earlier root of the analytic branch.

Neither method contains the other, and Best's published Table~1 settles it in
both directions. His table lists fourteen assumptions against some fifty
replacements, so a comparison needs a rule; ours is every (original,
replacement) pair he records ten or more times, which is \bestTop{} of them.
\bestMissing{} have no counterpart in our substitution table at all: they relax
an assumption to a bare operation --- an order to its \(<\), a module to its
scalar action --- and a table of named structures cannot name a target that is a
fragment of a class. Conversely, a pass that reads one proof term never
re-elaborates the weakened statement, and \incoherentPercent{}\% of our attempts
fail precisely there: a sibling assumption demands the class back, nothing
typechecks, and there is no term to inspect. Where both look at the same edge
they agree closely --- field to division ring is \bestFieldDR{} for him and
\ourFieldDR{} for us. The per-theorem overlap stays unmeasured, because his
metaprogram targets the previous prover and the previous library.

Gandhi, Tadipatri and Gowers \cite{gandhi2025} give a Lean tactic that
generalises the \emph{constants} in a proof, and set our target aside in one
sentence: ``it is unlikely that any proof-generalization algorithm would be
useful when applied to results in Lean's mathlib, since proofs there are already
designed to be as general as possible.'' We are careful about what our data does
to that. It does not show a generalisation algorithm is \emph{useful} on
mathlib; usefulness needs somebody to want the output, and nothing here has been
offered upstream. What it bears on is the stated reason. Proofs there are not
quite already as general as possible, at a measured rate of \ourYield{}\% of
attempts --- small, and not zero.

Baanen \cite{baanen2022} is the standing account of instance parameters in this
library, which is the object we manipulate. Mathlib's own maintenance report
\cite{baanen2025} describes an \texttt{unusedVariable} linter that already
catches hypotheses which turn out to be superfluous; anything a shipped linter
finds is not ours to report, which is one more reason the prior-art gate runs
before anything is counted.

\section{What this does not show}\label{sec:limits}

\paragraph{No novelty is claimed.} A reader will grant each survivor on sight.
The library's own search cannot prove them, which is a much weaker statement
than novelty, and it is the only one we make.

\paragraph{Nobody has asked for these.} None has been offered upstream, and the
difference between a weakening that is true and one that is wanted is the whole
value of the result. We have one piece of evidence on the second and it is not
ours: within sixteen days of our revision a Mathlib maintainer removed an unused
section variable in one of the files we patched, by hand, weakening it further
than we had. The \emph{kind} of change is wanted and gets merged; that is not
the same as anyone wanting ours. What we can offer is \patchKept{} one-line
patches a maintainer can check in a minute each (Section~\ref{sec:patches}) and
the evidence of Section~\ref{sec:used} that these theorems are used rather than
dead.

\paragraph{Nothing is claimed about how far the gap goes.} Each row records the
weakest class it is \emph{verified} to prove, and a search below that reached
lower classes we could not recompile; the published data carries those in a
column marked unverified, and no statistic here uses it. The ordinal measuring
the distance is also undefined for a pair of classes neither table places, so an
area with more unplaced classes looks shallower whether or not it is --- $8$ of
Order's $44$ survivors against $22$ of the other $299$. Depth is reported per
row and not aggregated. The invertibility results do not use it.

\paragraph{The pipeline has shipped defects, and they are documented.} Every one
had the same shape: something that resolved nothing looked exactly like
something that contained nothing, and both print zero. The artifact
\cite{artifact} carries the data, the tooling, a verifier that recompiles every
published row from a clean checkout, an \texttt{ENGINEERING.md} giving each
defect with its diagnosis, and one regression test per defect in
\texttt{tests/}.

\paragraph{The examined set is a sample, and not a random one.}
Areas were chosen rather than drawn, and within an area a theorem enters only
if the collector can locate and rewrite one of its binders. The composition is
therefore skewed: RingTheory is $2.2\times$ over-represented against its share
of the library and Order $0.7\times$ under-represented. What matters is whether
the skew moves the answer. Re-weighting each area's own rate by its share of the
corpus rather than by its share of what we examined gives \reweightedRate{}\%
across the \reweightAreas{} areas with at least 200 theorems examined, against
\theoremRate{}\% as measured --- a factor of $1.07$, and upward. The measured
rate is if anything conservative, but every figure here should be read as a
sample statistic and not a census.

\paragraph{The rate over the corpus is a floor, not an estimate.}
\mathlibWeakenablePercent{}\% divides \theoremsWeakened{} by every declaration
in the library, including those the search never looked at, and it is reported
that way because the alternative is to quietly choose a smaller denominator.
The library divides into three groups: \noBinderPercent{}\% of declarations
carry no instance binder in scope, so there is nothing to weaken and they are
out of scope by definition; \noTableClassPercent{}\% carry binders whose classes
neither of our tables covers, which is a limit of the tables; and
\reachablePercent{}\% carry a class we could weaken. We examined
\theoremsExamined{}, well under that \reachablePercent{}\%, because the
hand-table sweep ran \allAreasWord{} areas rather than the whole library and
because of the binder-locating limit above.

So \theoremRate{}\% is a measurement over the theorems the method examined, and
\mathlibWeakenablePercent{}\% is what that implies for the library only if the
unexamined behave like the examined. We have not tested that, and the second
number should be read as a lower bound rather than a rate.

\paragraph{One library, one revision.} Every row is relative to a named Mathlib
commit. The artifact publishes \(645\) rows --- the union of both weakening
tables, so more than the \allSurvivors{} this paper computes over --- and
\(644\) of them recompile from a clean checkout at that revision. The one that
does not is published too, with its error, because a dropped row is a claim
nobody can check.

\section{Open questions}\label{sec:questions}

We are not expert in these areas and offer the following as observations rather
than claims.

\begin{itemize}
\item Observation~\ref{obs:ultrametric} is an ultrametric bound over a
  commutative ring where the library states it over a field. Our descent stage
  reached a semiring, which would be the more interesting object, but we could
  not reproduce that compile and do not publish it. Is the bound over a
  commutative ring a studied object, and does it in fact go down to a semiring?
\item Observation~\ref{obs:withtop} admits exactly the extended non-negative
  structures. Are there results about $\mathbb{R}_{\ge 0} \cup \{\infty\}$
  currently stated over groups, so that the group is a wrapper around an
  argument that never used it?
\item Our hand table names no categorical class. Swept with the derived table,
  CategoryTheory gave \catSurvivors{} survivor from \catAttempts{} attempts, and
  \catContextPercent{}\% of attempts could not be reconstructed at all: a
  categorical hypothesis is usually a relation between objects, and this method
  replaces one assumption on one variable, so the question is less unanswered
than unaskable by this instrument.
\end{itemize}

\appendix

\section{Four worked examples}\label{app:obs}

Chosen by reading, not by any filter. All \allSurvivors{} survivors are in the
companion catalogue, grouped by area and weakening, each with the signature Lean
prints for it.

\begin{observation}[An ultrametric bound needs no field]\label{obs:ultrametric}
Let $R$ be a commutative ring and $v \colon R \to \mathbb{R}$ an absolute value with
$v(x+y) \le \max\{v(x),v(y)\}$. For finite $I,J$, $A = (a_{ji}) \in R^{J\times
I}$ and $x \in R^{I}$,
\[
  \sup_{j}\, v\!\Big(\sum_{i} a_{ji} x_i\Big)
  \;\le\; \Big(\sup_{j,i} v(a_{ji})\Big)\cdot \sup_{i} v(x_i),
\]
and with a factor $|I|$ on the right, the same holds without the ultrametric
hypothesis. Both are stated in the library over a field.
\end{observation}

\begin{observation}[Comparison after collapsing a top element]\label{obs:withtop}
Let $\alpha$ carry a partial order and a commutative associative addition with
a zero --- no subtraction ---
and let $u \colon \alpha \cup \{\top\} \to \alpha$ fix $\alpha$ and send $\top$
to $0$. Then $u$ is monotone away from the top, reflects order under the
expected side condition, and is non-negative when $\alpha$ is. Stated over an
ordered abelian group.
\end{observation}

\begin{observation}[A rational floor certificate needs no inverses]
If a rational certificate for an element of an ordered ring already carries the
invertibility datum it needs, the resulting statement about its floor holds over
the ring. It is stated over a field, and a decision procedure in the library
constructs the field instance for no other purpose than to apply it.
\end{observation}

\begin{observation}[Sup and inf are uniformly continuous without division]
Let $\alpha$ carry a linear order and a ring structure whose order is strict,
and nothing further --- in particular no inverses. For $\varepsilon$ and
$a_1,a_2,b_1,b_2$ in $\alpha$, if $\lvert a_1-b_1\rvert<\varepsilon$ and
$\lvert a_2-b_2\rvert<\varepsilon$ then
$\lvert a_1\sqcup a_2-b_1\sqcup b_2\rvert<\varepsilon$, and the same with
$\sqcap$. Both are stated in the library over a field. Each proof is one line,
and cites a lattice inequality and transitivity.
\end{observation}

\small
\begin{thebibliography}{9}
\setlength{\itemsep}{0pt}\setlength{\parsep}{0pt}\setlength{\parskip}{0pt}
\bibitem{best2021} Alex J. Best. \emph{Automatically Generalizing Theorems Using
Typeclasses}. Workshop on Formal Mathematics for Mathematicians, CICM 2021.
EasyChair preprint 6216.
\bibitem{alama2013} Jesse Alama. \emph{Eliciting Implicit Assumptions of Mizar
Proofs by Property Omission}. Journal of Automated Reasoning 50(2):123--133,
2013.
\bibitem{gandhi2025} Anshula Gandhi, Anand Rao Tadipatri and Timothy Gowers.
\emph{Automatically Generalizing Proofs and Statements}. ITP 2025, LIPIcs
vol.~352, article 12. \texttt{doi:10.4230/LIPIcs.ITP.2025.12}.

\bibitem{baanen2022} Anne Baanen. \emph{Use and Abuse of Instance Parameters in
the Lean Mathematical Library}. ITP 2022, LIPIcs vol.~237, article 4.

\bibitem{baanen2025} Anne Baanen et al. \emph{Growing Mathlib: maintenance of a
large scale mathematical library}. arXiv:2508.21593, 2025.

\bibitem{tauceti} The Tau Ceti Project. \emph{Tau Ceti: a repository of formal
mathematics implemented by AI contributors}. Swept at revision
\texttt{675a2d2c}. \url{https://github.com/TauCetiProject/TauCeti}

\bibitem{artifact} Data, tooling and a verifier for every row in this note.
\url{https://github.com/carlok/unused-assumptions}

\bibitem{companion} The \texttt{unstated-conclusions} project, 2026. Proof-term
roots and citation counts over the rows of this note; both are published row by
row in \cite{artifact}. \url{https://github.com/carlok/unstated-conclusions}

\bibitem{pons2002} Olivier Pons. \emph{Generalization in Type Theory Based Proof
Assistants}. Types for Proofs and Programs, LNCS, 2002.
\end{thebibliography}

\end{document}
